% Extended abstract for SE 2027 (Software Engineering, GI), two pages in the LNI format, on the paper
% "Trace Equation Systems for Program Verification" (doc/array-acceleration/teq.tex).
% Build: latexmk -pdf teq-abstract.tex
\documentclass[english]{lni}

\usepackage{mathtools}
\usepackage{listings}

\lstset{basicstyle=\ttfamily\small,columns=fullflexible,keepspaces=true}

\newcommand{\code}[1]{\mathsf{#1}}
\newcommand{\butlast}[1]{\lfloor #1 \rfloor}
\newcommand{\TRACE}[4]{\code{trace}(#1,\ \code{for}\ #2\ \code{do}\ #3 \mapsto #4)}
\newcommand{\MAPS}[3]{\code{maps}(\code{for}\ #1\ \code{yield}\ #2 \mapsto #3)}
\newcommand{\INVS}[3]{\code{invs}(\code{for}\ #1\ \code{check}\ #2 \mapsto #3)}
\newcommand{\rewrites}{\rightsquigarrow}

\addbibresource{references.bib}

\begin{document}

% TODO: the authors and affiliations of the paper; the venue of the paper once it is published
\title[Trace Equation Systems]{Trace Equation Systems for Program Verification}
\author[1]{Gidon Ernst}{gidon.ernst@sosy.ifi.lmu.de}{}
\affil[1]{LMU Munich\\Germany}
\editor{SE 2027}
\booktitle{Software Engineering 2027}
\yearofpublication{2027}
\maketitle

\begin{abstract}
Trace Equation Systems (TEQs) represent loops by terms for their computations and prove them by rewriting with an
extensible database of lemmas instead of inferring invariants. We summarize the main ideas of~\cite{teq}.
\end{abstract}


\section{Two Challenges of Loop Verification}

Proofs about loops rest on inductive invariants in the style of Floyd and Hoare, which poses two challenges. The
first is a \emph{representation gap}: the invariant is a formula over the states at the loop head that has to capture
what the loop has computed so far, e.g., a sum over a prefix of an array, but the logic cannot speak about the
computation itself. The second is the \emph{reliance on inference}: automatic tools discover even the most common
invariants anew for each program, which is brittle and rarely reaches aggregates or quantified properties of arrays.
Yet for recurring fragments, such as counters, array copies, or sums, the invariants are well known and could be
looked up instead.

\section{Trace Equation Systems}

The first idea is to make computations first-class terms. A loop is projected onto the \emph{traces} of the values of
its variables at the loop head, built by $\TRACE{x^0}{\butlast{\overline{Y}}}{x, \overline{y}}{u}$, a fold that keeps
all intermediate values: it starts with $x^0$ and applies the update $u$ once per element of the dependencies
$\butlast{\overline{Y}}$, e.g., the traces of other variables without their last value; $\code{current}$ selects the
last value. Each variable gets a trace of its own, which makes reasoning compositional; a counter $K$ anchors all
traces of a loop to its $n$ iterations. TEQs are generated by symbolic execution, with a case per path through the
body and an \emph{iteration constraint} ($\code{invs}$) with the past loop tests; the translation is sound and
relatively complete, and the TEQ is an invariant by construction.

\paragraph{Example.} The loop below copies an array $\code{b}$ into $\code{a}$; on the right is its TEQ:

\noindent
\begin{minipage}[c]{0.36\textwidth}
\begin{lstlisting}
for (i = 0; i < len; i++)
  a[i] = b[i];
assume(0 <= j && j < len);
assert(a[j] == b[j]);
\end{lstlisting}
\end{minipage}
\hfill
\begin{minipage}[c]{0.62\textwidth}
$\begin{aligned}
K &= \code{counter}(0, 1, n) \qquad (0 \le n) \\
I &= \TRACE{0}{\butlast{K}}{i, k}{i + 1} \\
A &= \TRACE{a^0}{\butlast{K}, \butlast{I}}{a, k, i}{a[i \coloneqq \code{b}[i]]} \\
& \INVS{\butlast{K}, \butlast{I}}{k, i}{i < \code{len}}
\end{aligned}$
\end{minipage}

\smallskip
\noindent
The goal is $\code{current}(A)[j] = \code{b}[j]$ at the exit, where $\neg(\code{current}(I) < \code{len})$.

\section{Proofs by a Lemma Database}

The second idea is to name known fragments. \emph{Introduction rules} replace a fragment by a named function, e.g., a
trace that adds a constant by a counter, a trace of stores by $\code{stores}$, a map of reads by $\code{selects}$.
\emph{Acceleration rules} relate these functions to the observations of the goal, e.g., the read of an array after the
stores of a loop. For the example:
\begin{align*}
& I \rewrites \code{counter}(0, 1, n),\;
  \code{current}(I) \rewrites n,\;
  \INVS{\butlast{I}}{i}{i < \code{len}} \rewrites n \le \code{len}, \\[2pt]
& \code{current}(A)[j]
  \rewrites \code{current}(\code{stores}(a^0, \butlast{I}, \MAPS{\butlast{I}}{i}{\code{b}[i]}))[j] \\
& \rewrites \code{current}(\code{stores}(a^0, \butlast{I}, \code{selects}(\code{b}, \butlast{I})))[j] \\
& \rewrites \code{earlier}(j, \code{selects}(\code{b}, \butlast{\code{counter}(0, 1, n)}))
  \rewrites \code{b}[j]
\end{align*}
The loop test and its negation at the exit give $n = \code{len}$, so $0 \le j < n$, and the goal is
$\code{b}[j] = \code{b}[j]$. Nothing is guessed: the map of the values stored by the loop is synthesized while the
lemma for $\code{stores}$ is matched, and quantifiers are instantiated by matching.

The lemmas are generic in the sequences of keys and values. Reasoning is incremental, like with the simplifier of a
proof assistant; open subgoals point to missing lemmas. The database is extensible by new concepts, such as sums or
maxima, and we envision it as shared across tools; its lemmas are proved once, in Lean, not once per program.

\section{Implementation and Results}

TEQs are implemented in the toolchain Cuvée, with a deduction engine for the rewriting (congruence closure, a
simplex, matching modulo definitions) and 153 lemmas, all proved in Lean. On the safe tasks of SV-COMP's
\texttt{Arrays}, it proves more than any participant of SV-COMP 2026 (148 against 68), and 88 of the 125 aggregate
properties of MonoCera's benchmarks (against 41), without a wrong answer. Tools that infer invariants are stronger on
counterexamples and numeric properties; TEQs complement them where knowledge about a computation is written down once
and reused.

\printbibliography

\end{document}
